For each $N\in\mathbb N$, define
\[
\delta(N):=\min\left\{\left|1-\sum_{n\in A}\frac1n\right|:\;
A\subseteq\{1,2,\dots,N\},\;
\forall S\subseteq A,\ \sum_{n\in S}\frac1n\neq 1
\right\}.
\]
Equivalently, we look for a subset $A$ whose reciprocal sum is as close as possible to $1$,
subject to the constraint that \emph{no} subsubset $S\subseteq A$ has reciprocal sum exactly $1$.

The question asks for the asymptotic behavior of $\delta(N)$ as $N\to\infty$; in particular,
whether
\[
\delta(N)=e^{-(c+o(1))N}
\quad\text{for some constant }c\in(0,1).
\]

\textbf{Stress points / ambiguity checks.}
\begin{itemize}
\item The constraint is on \emph{all} subsubsets $S\subseteq A$, including $S=A$; thus any $A$ with $\sum_{n\in A}1/n=1$ is forbidden.
\item Without the constraint, the minimal nonzero distance to $1$ is at least $1/\mathrm{lcm}(1,\dots,N)$ by a denominator argument (proved below).
\item The statement ``$\delta(N)=e^{-(c+o(1))N}$'' should be read as a conjectural asymptotic, not as already-established.
\end{itemize}
